\documentclass[../../zusammenfassung_AN.tex]{subfiles}
\begin{document}
	
	% -------------------------------------------------------
	\subsubsection{Folgen}
	% -------------------------------------------------------
	
	\begin{tcolorbox}[formelbox]
		Eine \textbf{Folge} ist eine Abbildung $n \mapsto a_n$:
		\[
		(a_n)_{n\in\mathbb{N}}, \quad a_n \in \mathbb{R}.
		\]
		Die Folge \textbf{konvergiert} gegen $L \in \mathbb{R}$, falls
		$\displaystyle\lim_{n\to\infty} a_n = L$.
	\end{tcolorbox}
	
	\paragraph{Wichtige Folgentypen}
	
	\begin{tcolorbox}[dglbox, title={Arithmetische Folge \quad $a_n = a_1 + (n-1)d$}]
		\[
		\lim_{n\to\infty} a_n =
		\begin{cases}
			+\infty & d>0 \\
			-\infty & d<0 \\
			a_1     & d=0 \quad \text{(einziger konvergenter Fall)}
		\end{cases}
		\]
	\end{tcolorbox}
	
	\begin{tcolorbox}[dglbox, title={Geometrische Folge \quad $a_n = a_1\,q^{n-1}$}]
		\[
		\lim_{n\to\infty} a_n =
		\begin{cases}
			0                        & |q|<1 \\
			a_1                      & q=1 \\
			\text{existiert nicht}   & q=-1 \\
			\pm\infty                & |q|>1
		\end{cases}
		\]
	\end{tcolorbox}
	
	\begin{tcolorbox}[formelbox]
		\textbf{Potenzfolgen} $a_n=n^p$: \quad
		$p>0$: divergiert $(\to\infty)$ \quad
		$p=0$: konstant \quad
		$p<0$: $\to 0$
		
		\medskip
		\textbf{Quotienten von Polynomen} $a_n = \dfrac{n^k}{n^m}$: Vergleich der höchsten Potenzen:
		\[
		k < m \;\Rightarrow\; 0 \qquad
		k = m \;\Rightarrow\; \text{Quotient der Leitkoeffizienten} \qquad
		k > m \;\Rightarrow\; \text{divergent}
		\]
	\end{tcolorbox}
	
	% -------------------------------------------------------
	\subsubsection{Reihen}
	% -------------------------------------------------------
	
	\begin{tcolorbox}[formelbox]
		Eine \textbf{Reihe} ist die Summe einer Folge:
		$\displaystyle\sum_{n=1}^{\infty} a_n$
		
		\medskip
		\textbf{Notwendige Bedingung für Konvergenz:}
		\[
		\sum_{n=1}^{\infty} a_n \text{ konvergiert} \;\Rightarrow\; a_n \xrightarrow{n\to\infty} 0
		\]
		Gilt $a_n \not\to 0$, so ist die Reihe \textbf{sicher divergent}.
	\end{tcolorbox}
	
	\paragraph{Arithmetische Reihe (endlich)}
	
	\begin{tcolorbox}[formelbox]
		\[
		S_n = \sum_{k=1}^{n} a_k = \frac{n}{2}\bigl(2a_1+(n-1)d\bigr) = \frac{n}{2}(a_1+a_n)
		\]
		Die unendliche arithmetische Reihe \textbf{divergiert} immer (ausser $a_1=d=0$). \\
		\textit{(Gausstrick: Erste und letzte Glieder paarweise addieren)}
	\end{tcolorbox}
	
	\paragraph{Geometrische Reihe}
	
	\begin{tcolorbox}[dglbox, title={Geometrische Reihe \quad $\sum_{n=0}^{\infty} a\,q^n$}]
		\textbf{Konvergenz:} $|q|<1$ \qquad
		\textbf{Summe:}
		\[
		S = \frac{a}{1-q}
		\]
	\end{tcolorbox}
	
	\paragraph{Harmonische Reihe}
	
	\begin{tcolorbox}[hinweis]
		\[
		\sum_{n=1}^{\infty} \frac{1}{n} \quad \textbf{divergiert}.
		\]
	\end{tcolorbox}
	
	\paragraph{$p$-Reihe}
	
	\begin{tcolorbox}[formelbox]
		\[
		\sum_{n=1}^{\infty} \frac{1}{n^p}:\qquad
		p>1 \;\Rightarrow\; \text{konvergent} \qquad
		p\le 1 \;\Rightarrow\; \text{divergent}
		\]
	\end{tcolorbox}
	
	\paragraph{Potenzreihen (endlich, Beispiel $\sum k^2$)}
	
	\begin{tcolorbox}[formelbox]
		\[
		\sum_{k=1}^{n} k^2 = \frac{n(n+1)(2n+1)}{6} \qquad
		\sum_{k=1}^{n} k   = \frac{n(n+1)}{2} \qquad
		\sum_{k=1}^{n} k^3 = \left(\frac{n(n+1)}{2}\right)^2
		\]
	\end{tcolorbox}
	
	\begin{tcolorbox}[hinweis]
		\textbf{Bernoulli-Trick} (zur Herleitung von $\sum k^p$): Betrachte die Teleskopsumme
		\[
		\sum_{k=1}^{n}\bigl((k+1)^{p+1}-k^{p+1}\bigr) = (n+1)^{p+1}-1.
		\]
		Linke Seite mit Binomischem Lehrsatz ausmultiplizieren, bekannte Summenformeln einsetzen und nach $\sum k^p$ auflösen.
	\end{tcolorbox}
	
	\paragraph{Alternierende Reihen}
	
	\begin{tcolorbox}[dglbox, title={Leibniz-Kriterium}]
		$\sum(-1)^n a_n$ konvergiert, falls:
		\begin{enumerate}[noitemsep, topsep=2pt]
			\item $a_n \ge 0$
			\item $(a_n)$ monoton fallend
			\item $a_n \to 0$
		\end{enumerate}
	\end{tcolorbox}
	
	\paragraph{Teleskopsummen}
	
	\begin{tcolorbox}[formelbox]
		\[
		\sum_{k=1}^{n}\bigl(f(k+1)-f(k)\bigr) = f(n+1)-f(1)
		\]
	\end{tcolorbox}
	
	% -------------------------------------------------------
	\subsubsection{Konvergenzkriterien}
	% -------------------------------------------------------
	
	\begin{tcolorbox}[dglbox, title={Konvergenzkriterien im Überblick}]
		\textbf{Majorantenkriterium:}
		\[
		0 \le a_n \le b_n \;\text{ und }\; \sum b_n \text{ konvergiert}
		\;\Rightarrow\; \sum a_n \text{ konvergiert}
		\]
		\textbf{Minorantenkriterium:}
		\[
		a_n \ge b_n \ge 0 \;\text{ und }\; \sum b_n \text{ divergiert}
		\;\Rightarrow\; \sum a_n \text{ divergiert}
		\]
		\textbf{Quotientenkriterium:}
		\[
		L = \lim_{n\to\infty}\left|\frac{a_{n+1}}{a_n}\right|:\qquad
		L<1 \;\Rightarrow\; \text{konvergent} \qquad
		L>1 \;\Rightarrow\; \text{divergent} \qquad
		L=1 \;\Rightarrow\; \text{keine Aussage}
		\]
	\end{tcolorbox}
	
	\begin{tcolorbox}[hinweis]
		\textbf{Prüfungsstrategie:}
		\begin{enumerate}[noitemsep]
			\item Gilt $a_n \to 0$? (falls nein: divergent)
			\item Geometrische Reihe erkennen?
			\item Vergleich mit $p$-Reihe (Majorante/Minorante)
			\item Alternierend $\Rightarrow$ Leibniz
			\item Quotientenkriterium anwenden
		\end{enumerate}
	\end{tcolorbox}
	
\end{document}